

% ============================================================
%  CAPÍTULO: INTEGRALES IMPROPIAS
%  Libro: Notas de Cálculo — anfalearNotasCalculo
%  Requiere en el preámbulo del libro:
%    \usepackage{pgfplots}
%    \pgfplotsset{compat=1.18}
%    \usepgfplotslibrary{fillbetween}
%    \usetikzlibrary{patterns,arrows.meta}
%    \usepackage{enumerate}   % para [(a)]
% ============================================================

\chapter{Integrales Impropias}


En el marco del cálculo de Riemann, la integral definida $\int_{a}^{b} f(x)\,dx$ se construye
sobre intervalos cerrados y acotados $[a, b]$. Sin embargo, el hilo conductor que conecta
la geometría con el análisis nos invita a extender esta noción: si el área bajo una curva es la
suma de infinitos rectángulos infinitesimales, ¿qué ocurre cuando la región de integración se
prolonga indefinidamente hacia la derecha o hacia la izquierda? Para dar sentido matemático a
fenómenos físicos, modelos probabilísticos y problemas de ingeniería donde el dominio no está
acotado, introducimos el concepto de \textbf{integral impropia de primer tipo}.

\begin{definition}[Integral impropia con límite infinito]
Sea $f(x)$ una función continua en $[a, \infty)$. Definimos la integral impropia con límite
superior infinito como el límite de integrales propias:
\[
  \int_{a}^{\infty} f(x)\,dx = \lim_{b \to \infty} \int_{a}^{b} f(x)\,dx.
\]
De manera análoga, si el límite inferior es infinito:
\[
  \int_{-\infty}^{b} f(x)\,dx = \lim_{a \to -\infty} \int_{a}^{b} f(x)\,dx.
\]
Si el intervalo es $(-\infty,\infty)$, la integral se define separando en un punto arbitrario
$c$ (frecuentemente $c=0$):
\[
  \int_{-\infty}^{\infty} f(x)\,dx
  = \int_{-\infty}^{c} f(x)\,dx + \int_{c}^{\infty} f(x)\,dx.
\]
\end{definition}

\begin{theorem}[Criterio de convergencia]
Las integrales impropias de límite infinito \textbf{convergen} si y solo si los límites
involucrados en su definición existen y son números reales finitos. Si el límite es $\pm\infty$
o no existe, se dice que la integral \textbf{diverge}.
\begin{proof}
La integral propia $\int_{a}^{b} f(x)\,dx$ define una función continua de $b$. El límite
$\lim_{b\to\infty}$ existe si y solo si la sucesión de sumas de Riemann parciales está acotada y
tiende a un valor finito. La divergencia ocurre cuando el crecimiento acumulado del área supera
cualquier cota finita, o cuando hay oscilaciones no amortiguadas que impiden la existencia del
límite.
\end{proof}
\end{theorem}

\begin{rem}Es importante tener en cuenta algunas observaciones básicas:
\begin{itemize}
    \item \textbf{Jerarquía de crecimiento:} Las funciones exponenciales con exponente negativo
          ($e^{-kx}$, $k>0$) decaen mucho más rápido que cualquier función racional $1/x^p$. Esto
          garantiza la convergencia de integrales como $\int_0^\infty x^n e^{-x}\,dx$.

    \item \textbf{Criterio $p$ para límites infinitos:}
          $\displaystyle\int_{1}^{\infty}\frac{dx}{x^p}$ converge
          $\iff p>1$.

    \item \textbf{Error frecuente:} No evalúe directamente $F(\infty)-F(a)$. Siempre plantee el
          límite formal. La expresión $\infty - \infty$ es una \emph{forma indeterminada} que
          requiere tratamiento analítico riguroso; en particular, aplique la regla de L'Hôpital o
          la jerarquía de crecimiento antes de concluir.
\end{itemize}
\end{rem}

A continuación ilustramos el procedimiento paso a paso, desde el caso más sencillo hasta el
que involucra límites en ambos extremos.

% ------------------------------------------------------------
\begin{example}[Área bajo una campana exponencial decreciente]
Calcule el valor de la integral impropia
$\displaystyle\int_{0}^{\infty} x\,e^{-x}\,dx$.
\end{example}
\begin{myproof}
\textbf{Paso 1: Análisis gráfico y región.}

La función $f(x)=xe^{-x}$ es positiva para $x>0$, parte del origen, alcanza su máximo en
$x=1$ (pues $f'(x)=e^{-x}(1-x)=0\Rightarrow x=1$, con $f(1)=e^{-1}$) y decae asintóticamente a
$0$ cuando $x\to\infty$. La región de integración es el área infinita bajo la curva desde $x=0$
hacia la derecha; visualizamos el límite superior $b$ como una \textbf{frontera móvil} que avanza
indefinidamente.

\begin{center}
\begin{tikzpicture}
\begin{axis}[
    width=0.62\textwidth, height=0.38\textheight,
    axis lines=middle,
    xlabel={$x$}, ylabel={$y$},
    xmin=-0.4, xmax=6.3, ymin=-0.04, ymax=0.44,
    domain=0:6, samples=150,
    xtick={0,1,2,3,4,5},
    ytick={0,0.1,0.2,0.3},
    legend pos=north east
]
\addplot[name path=f, thick, blue!80!black] {x*exp(-x)};
\addlegendentry{$y=xe^{-x}$}
\path[name path=eje] (axis cs:0,0) -- (axis cs:6,0);
\addplot[blue!20, opacity=0.85]
    fill between[of=f and eje, soft clip={domain=0:4.6}];
\draw[dashed, red!70!black, thick]
    (axis cs:4.6,0) node[below] {$b$}
    -- (axis cs:4.6,{4.6*exp(-4.6)});
\filldraw[red!80!black]
    (axis cs:1,{exp(-1)}) circle (1.8pt)
    node[above right] {\small Máx.\ $\bigl(1,\,e^{-1}\bigr)$};
\node[blue!70!black] at (axis cs:2.2,0.16) {Región $[0,b]$};
\end{axis}
\end{tikzpicture}
\end{center}

\textbf{Paso 2: Elemento diferencial.}
Considerando franjas verticales de altura $f(x)$ y base $dx$:
\[
  dA = xe^{-x}\,dx.
\]

\textbf{Paso 3: Límites de integración.}
El intervalo de integración es $[0,b]$; se toma el límite $b\to\infty$ para barrer toda la región.

\textbf{Paso 4: Evaluación.}
Aplicamos integración por partes con $u=x$ y $dv=e^{-x}\,dx$, lo que da $du=dx$ y $v=-e^{-x}$:
\begin{align*}
\int_{0}^{b} xe^{-x}\,dx
  &= \Bigl[-xe^{-x}\Bigr]_0^b + \int_0^b e^{-x}\,dx \\
  &= \bigl(-be^{-b}+0\bigr) + \Bigl[-e^{-x}\Bigr]_0^b \\
  &= -be^{-b} - e^{-b} + 1
   = 1 - e^{-b}(b+1).
\end{align*}
Tomando el límite (la forma $\infty/\infty$ se resuelve con L'Hôpital):
\begin{align*}
\lim_{b\to\infty}\bigl[1-e^{-b}(b+1)\bigr]
  &= 1 - \lim_{b\to\infty}\frac{b+1}{e^{b}}
   \;\stackrel{\text{L'H}}{=}\;
   1 - \lim_{b\to\infty}\frac{1}{e^{b}} = 1 - 0 = 1.
\end{align*}
\[
  \boxed{\int_{0}^{\infty} xe^{-x}\,dx = 1.}
\]
\end{myproof}

% ------------------------------------------------------------
\begin{example}[Integral con límites infinitos en ambos extremos]
Determine si $\displaystyle\int_{-\infty}^{\infty}\frac{dx}{1+x^2}$ converge y, en caso
afirmativo, calcule su valor.
\end{example}
\begin{myproof}
\textbf{Paso 1: Análisis gráfico y región.}

La función $f(x)=\tfrac{1}{1+x^2}$ es \textbf{par}, positiva en todo $\mathbb{R}$, alcanza su
máximo $f(0)=1$ y posee asíntota horizontal $y=0$. La región bajo la curva se extiende
infinitamente a ambos lados del eje $y$.

\begin{center}
\begin{tikzpicture}
\begin{axis}[
    width=0.62\textwidth, height=0.38\textheight,
    axis lines=middle,
    xlabel={$x$}, ylabel={$y$},
    xmin=-6.2, xmax=6.2, ymin=-0.08, ymax=1.18,
    domain=-6:6, samples=150,
    xtick={-5,-3,-1,0,1,3,5},
    ytick={0,0.5,1},
    legend pos=north east
]
\addplot[name path=f, thick, purple!80!black] {1/(1+x^2)};
\addlegendentry{$y=\tfrac{1}{1+x^2}$}
\path[name path=eje] (axis cs:-6,0) -- (axis cs:6,0);
\addplot[purple!18, opacity=0.90]
    fill between[of=f and eje, soft clip={domain=-5.5:5.5}];
\filldraw[black] (axis cs:0,1) circle (1.8pt) node[right] {\small $(0,1)$};
\end{axis}
\end{tikzpicture}
\end{center}

\textbf{Paso 2: Elemento diferencial.}
\[
  dA = \frac{dx}{1+x^2}.
\]

\textbf{Paso 3: Límites de integración.}
Separamos en $c=0$ para manejar los dos infinitos de forma \emph{independiente}:
\[
  \int_{-\infty}^{\infty}\frac{dx}{1+x^2}
  = \int_{-\infty}^{0}\frac{dx}{1+x^2}
  + \int_{0}^{\infty}\frac{dx}{1+x^2}.
\]

\textbf{Paso 4: Evaluación.}
Usando $\int\tfrac{dx}{1+x^2}=\arctan(x)+C$:
\begin{align*}
\int_{-\infty}^{\infty}\frac{dx}{1+x^2}
  &= \lim_{a\to -\infty}\Bigl[\arctan x\Bigr]_a^{0}
   + \lim_{b\to\infty}\Bigl[\arctan x\Bigr]_0^{b} \\
  &= \Bigl(0-\bigl(-\tfrac{\pi}{2}\bigr)\Bigr)
   + \Bigl(\tfrac{\pi}{2}-0\Bigr)
   = \frac{\pi}{2}+\frac{\pi}{2}.
\end{align*}
\[
  \boxed{\int_{-\infty}^{\infty}\frac{dx}{1+x^2} = \pi.}
\]
\end{myproof}

% ============================================================
\section{Funciones de densidad de probabilidad}
% ============================================================

El hilo conductor entre el cálculo integral y la estadística reside en la interpretación
geométrica del área: mientras que en geometría el área representa una magnitud espacial, en
probabilidad representa la \textbf{certeza total}. Las variables aleatorias continuas, como la vida
útil de un componente electrónico o la resistencia de un material, modelan sus distribuciones
mediante funciones no negativas cuya área total bajo la curva debe ser exactamente $1$. Toda
integral que aparezca en este marco es, en esencia, una integral impropia de primer tipo.

\begin{definition}[Función de densidad de probabilidad]
Una función $f(x)$ es una \textbf{función de densidad de probabilidad} (FDP) para una variable
aleatoria continua $X$ si satisface:
\begin{enumerate}
    \item $f(x)\geq 0$ para todo $x\in\mathbb{R}$.
    \item $\displaystyle\int_{-\infty}^{\infty} f(x)\,dx = 1$.
\end{enumerate}
La probabilidad de que $X$ tome valores en el intervalo $[a,b]$ es:
\[
  P(a\leq X\leq b) = \int_{a}^{b} f(x)\,dx.
\]
\end{definition}

\begin{rem} 
Antes de utilizar una función como FDP, realice siempre las siguientes comprobaciones:
\begin{enumerate}
 \item \textbf{no negatividad:}
          Compruebe analíticamente que $f(x)\geq 0$ en todo su dominio. Basta con inspeccionar el
          signo del numerador y del denominador.
    \item \textbf{normalización:}
          Calcule $\int_{-\infty}^{\infty}f(x)\,dx$. Si el resultado es $1$, la función es válida.
          Si es una constante $C\neq 1$, multiplique por $1/C$ para obtener la FDP normalizada.
    \item \textbf{Advertencia fundamental:}
          $f(x)$ \textbf{no} es una probabilidad puntual; puede tomar valores mayores que $1$.
          Solo la integral $\int_a^b f(x)\,dx$ tiene significado probabilístico. En particular,
          $P(X=c)=0$ para todo valor fijo $c$, pues la integral sobre un conjunto de medida cero
          es cero.
\end{enumerate}
   

\end{rem}

% -- Ejemplo adicional: normalización ----------------------------
\begin{example}[Determinación de la constante de normalización]
Encuentre el valor de la constante $k>0$ para que la función
\[
  f(x) = \begin{cases} ke^{-3x} & x\geq 0 \\ 0 & x<0 \end{cases}
\]
sea una función de densidad de probabilidad. Luego, dibuje la función resultante.
\end{example}
\begin{myproof}
\textbf{Paso 1: Análisis gráfico y región.}

Para cualquier $k>0$, la función $ke^{-3x}$ es positiva y estrictamente decreciente sobre
$[0,\infty)$, por lo que la \emph{Verificación 1} se cumple automáticamente. La condición de
normalización (Verificación 2) es la ecuación que determina $k$: exigimos que el área total bajo
la curva sea exactamente $1$.

\begin{center}
\begin{tikzpicture}
\begin{axis}[
    width=0.58\textwidth, height=0.36\textheight,
    axis lines=middle,
    xlabel={$x$}, ylabel={$f(x)$},
    xmin=-0.3, xmax=2.2, ymin=-0.12, ymax=3.4,
    domain=0:2, samples=120,
    xtick={0,0.5,1,1.5,2}, ytick={0,1,2,3}
]
\addplot[name path=f, thick, orange!85!red] {3*exp(-3*x)};
\addlegendentry{$f(x)=3e^{-3x}$}
\path[name path=eje] (axis cs:0,0) -- (axis cs:2,0);
\addplot[orange!22, opacity=0.9]
    fill between[of=f and eje, soft clip={domain=0:2}];
\filldraw[black] (axis cs:0,3) circle (1.8pt)
    node[right] {\small $(0,\,k)$};
\node[orange!70!red] at (axis cs:0.85,0.90) {Área$=1$};
\end{axis}
\end{tikzpicture}
\end{center}

\textbf{Paso 2: Elemento diferencial.}
\[
  dA = ke^{-3x}\,dx.
\]

\textbf{Paso 3: Límites de integración.}
La función es no nula únicamente sobre $[0,\infty)$, así que la condición de normalización es
$\int_0^{\infty}ke^{-3x}\,dx = 1$.

\textbf{Paso 4: Evaluación.}
\begin{align*}
\int_0^{\infty} ke^{-3x}\,dx
  &= k\lim_{b\to\infty}\int_0^b e^{-3x}\,dx
   = k\lim_{b\to\infty}\Bigl[-\tfrac{1}{3}e^{-3x}\Bigr]_0^b \\
  &= k\lim_{b\to\infty}\Bigl(-\tfrac{1}{3}e^{-3b}+\tfrac{1}{3}\Bigr)
   = k\cdot\frac{1}{3} = 1.
\end{align*}
Por tanto $k=3$, y la FDP normalizada es:
\[
  \boxed{f(x)=3e^{-3x},\quad x\geq 0.}
\]
\end{myproof}

\begin{theorem}[Media y varianza de una variable aleatoria continua]
La \textbf{media} $\mu$ (valor esperado) y la \textbf{varianza} $\sigma^2$ de una variable
aleatoria continua con FDP $f(x)$ se definen como:
\[
  \mu = E(X) = \int_{-\infty}^{\infty} x\,f(x)\,dx,
  \qquad
  \sigma^2 = V(X) = \int_{-\infty}^{\infty}(x-\mu)^2 f(x)\,dx.
\]
Existe además la \textbf{fórmula computacional}: $\sigma^2 = E(X^2)-[E(X)]^2$.
\begin{proof}
Aplicando la linealidad de la esperanza a $g(X)=(X-\mu)^2$ y expandiendo el cuadrado:
\begin{align*}
\sigma^2
  &= E\bigl[(X-\mu)^2\bigr]
   = E\bigl[X^2 - 2\mu X + \mu^2\bigr] \\
  &= E(X^2) - 2\mu\,E(X) + \mu^2
   = E(X^2) - 2\mu^2 + \mu^2
   = E(X^2) - \mu^2.
\end{align*}
La convergencia de estas integrales está garantizada si los momentos de orden 1 y 2 son finitos.
\end{proof}
\end{theorem}

\begin{rem}La Interpretación geométrica y estadística de estas integrales son:
\begin{itemize}
    \item \textbf{Centro de masa:} $\mu$ actúa como el punto de equilibrio (centro de gravedad)
          de la distribución. Si la FDP es simétrica respecto a $x=\mu$, la media coincide con el
          eje de simetría.
    \item \textbf{Dispersión:} $\sigma^2$ cuantifica la «anchura» de la distribución. Una
          $\sigma^2$ pequeña indica alta concentración alrededor de $\mu$; una grande indica alta
          dispersión.
    \item \textbf{Preferencia computacional:} Use siempre $\sigma^2=E(X^2)-\mu^2$ para cálculos
          algebraicos; evita desarrollar el cuadrado $(x-\mu)^2$ bajo el signo integral, lo que
          simplifica notablemente la integración por partes.
\end{itemize}
\end{rem}

% ------------------------------------------------------------
\begin{example}[Distribución exponencial: vida útil de un sensor]
La vida útil $X$ (en miles de horas) de un sensor de presión sigue la FDP
\[
  f(x) = \begin{cases} 2e^{-2x} & x\geq 0 \\ 0 & x<0. \end{cases}
\]
\begin{enumerate}[(a)]
  \item Calcule la probabilidad de que el sensor falle antes de las $500$ horas ($X<0.5$).
  \item Determine la vida útil media $\mu$ y la varianza $\sigma^2$.
\end{enumerate}
\end{example}
\begin{myproof}
\textbf{Paso 1: Análisis gráfico y región.}

La función parte de $f(0)=2$, decrece exponencialmente y tiende a $0$. La región azul representa
$P(X<0.5)$; la totalidad del área bajo la curva es $1$.

\begin{center}
\begin{tikzpicture}
\begin{axis}[
    width=0.60\textwidth, height=0.37\textheight,
    axis lines=middle,
    xlabel={$x$}, ylabel={$f(x)$},
    xmin=-0.2, xmax=2.5, ymin=-0.10, ymax=2.25,
    domain=0:2.3, samples=130,
    xtick={0,0.5,1,1.5,2}, ytick={0,0.5,1,1.5,2}
]
\addplot[name path=f, thick, blue!75!black] {2*exp(-2*x)};
\addlegendentry{$f(x)=2e^{-2x}$}
\path[name path=eje] (axis cs:0,0) -- (axis cs:2.3,0);
\addplot[blue!22, opacity=0.90]
    fill between[of=f and eje, soft clip={domain=0:0.5}];
\draw[dashed, red!70!black, thick]
    (axis cs:0.5,0) node[below] {$0.5$}
    -- (axis cs:0.5,{2*exp(-1)});
\node[blue!65!black] at (axis cs:0.22,0.5)
    {\small $P(X{<}0.5)$};
\end{axis}
\end{tikzpicture}
\end{center}

\textbf{Paso 2: Elemento diferencial.}
\[
  dP = f(x)\,dx = 2e^{-2x}\,dx.
\]

\textbf{Paso 3: Límites de integración.}
\begin{enumerate}[(a)]
  \item Probabilidad: intervalo $[0,\,0.5]$.
  \item Media: $\int_0^\infty x\cdot 2e^{-2x}\,dx$.\quad
        Segundo momento: $\int_0^\infty x^2\cdot 2e^{-2x}\,dx$.
\end{enumerate}

\textbf{Paso 4: Evaluación.}

\textit{(a) Probabilidad de falla temprana:}
\begin{align*}
P(X<0.5) &= \int_0^{0.5} 2e^{-2x}\,dx
           = \Bigl[-e^{-2x}\Bigr]_0^{0.5}
           = -e^{-1}+e^{0} = 1-e^{-1}.
\end{align*}

\textit{(b) Media} (por partes: $u=x$, $dv=2e^{-2x}\,dx \Rightarrow v=-e^{-2x}$):
\begin{align*}
\mu &= \int_0^{\infty} x\cdot 2e^{-2x}\,dx
     = \Bigl[-xe^{-2x}\Bigr]_0^{\infty}+\int_0^\infty e^{-2x}\,dx
     = 0 + \Bigl[-\tfrac{1}{2}e^{-2x}\Bigr]_0^\infty = \frac{1}{2}.
\end{align*}

\textit{Segundo momento} (por partes dos veces, con $u=x^2$, $dv=2e^{-2x}\,dx$):
\begin{align*}
E(X^2)
  &= \int_0^\infty x^2\cdot 2e^{-2x}\,dx
   = \Bigl[-x^2 e^{-2x}\Bigr]_0^\infty + \int_0^\infty 2xe^{-2x}\,dx
   = 0 + \frac{1}{2} = \frac{1}{2}.
\end{align*}

\textit{Varianza} (fórmula computacional):
\begin{align*}
\sigma^2 &= E(X^2)-\mu^2 = \frac{1}{2}-\frac{1}{4} = \frac{1}{4}.
\end{align*}
\[
  \boxed{P(X<0.5)=1-e^{-1}\approx 0.6321,\qquad
         \mu=\tfrac{1}{2}\text{ (mil horas)},\qquad
         \sigma^2=\tfrac{1}{4}.}
\]
\end{myproof}

% ============================================================
\section{Integrandos no acotados y discontinuidades}
% ============================================================

Hasta ahora hemos extendido los \emph{límites} de integración hacia el infinito. Ahora abordaremos
la segunda clase de integrales impropias: aquellas donde la función $f(x)$ \textbf{no está acotada}
dentro del intervalo de integración. \textbf{El hilo conductor} aquí es la ruptura del Teorema
Fundamental del Cálculo (TFC): el TFC requiere continuidad de $f$ en $[a,b]$. Si $f$ tiende a
infinito en un extremo o en un punto interior, las sumas de Riemann pueden diverger o requerir un
tratamiento por límites laterales.

\begin{definition}[Integral impropia de tipo II]
Sea $f$ continua en $[a,b)$ con $\lim_{x\to b^-}|f(x)|=\infty$. Definimos:
\[
  \int_{a}^{b} f(x)\,dx = \lim_{t\to b^-}\int_{a}^{t} f(x)\,dx.
\]
Si la discontinuidad está en el extremo izquierdo $x=a$:
\[
  \int_{a}^{b} f(x)\,dx = \lim_{t\to a^+}\int_{t}^{b} f(x)\,dx.
\]
\end{definition}

\begin{theorem}[Convergencia con discontinuidad interior]
Si $f$ tiene una discontinuidad infinita en $c\in(a,b)$, entonces:
\[
  \int_{a}^{b} f(x)\,dx = \int_{a}^{c} f(x)\,dx + \int_{c}^{b} f(x)\,dx,
\]
y la integral total converge si y solo si \textbf{ambas} integrales del lado derecho convergen.
\begin{proof}
Por la aditividad del intervalo de integración y la independencia de los límites laterales
$t\to c^-$ y $s\to c^+$. Si uno de ellos diverge, la suma no puede asignarse un valor finito,
destruyendo la definición de convergencia.
\end{proof}
\end{theorem}

\begin{rem} Al intentar resolver cualquier tipo de integral, debe seguirse el siguiente protocolo de verificación antes de aplicar el TFC
\begin{enumerate}
    \item  Antes de escribir $F(b)-F(a)$, verifique si $f(x)$ tiene
          \textbf{asíntotas verticales} en $[a,b]$. Busque los ceros del denominador y los puntos
          donde el exponente es negativo.
    \item \textbf{Criterio $p$ para integrandos infinitos:}
          $\displaystyle\int_0^1\frac{dx}{x^p}$ converge $\iff p<1$.
          \textbf{Note la inversión} respecto al criterio de límites infinitos.
    \item \textbf{Trampa clásica:}
          $\displaystyle\int_{-1}^{1}\frac{dx}{x^2}\neq -2$.
          La función es \emph{positiva}; obtener un resultado negativo o finito sin plantear el
          límite en $x=0$ indica aplicación incorrecta del TFC.
\end{enumerate}
\end{rem}

% ------------------------------------------------------------
\begin{example}[Discontinuidad en el extremo inferior]
Analice la convergencia de $\displaystyle\int_{0}^{1}\frac{dx}{\sqrt{x}}$.
\end{example}
\begin{myproof}
\textbf{Paso 1: Análisis gráfico y región.}

La función $f(x)=x^{-1/2}$ crece sin cota cuando $x\to 0^+$, pero decrece rápidamente. La región
es infinita en \emph{altura} pero confinada horizontalmente al intervalo $(0,1]$; la intuición
visual sugiere que el área podría ser finita porque la función cae muy rápido.

\begin{center}
\begin{tikzpicture}
\begin{axis}[
    width=0.56\textwidth, height=0.37\textheight,
    axis lines=middle,
    xlabel={$x$}, ylabel={$y$},
    xmin=-0.08, xmax=1.25, ymin=-0.18, ymax=4.1,
    restrict y to domain=0:4.1,
    domain=0.006:1.18, samples=180,
    xtick={0,0.25,0.5,0.75,1}, ytick={0,1,2,3,4}
]
\addplot[name path=f, thick, teal!80!black] {1/sqrt(x)};
\addlegendentry{$y=x^{-1/2}$}
\path[name path=eje] (axis cs:0,0) -- (axis cs:1.18,0);
\addplot[teal!22, opacity=0.90]
    fill between[of=f and eje, soft clip={domain=0.006:1}];
\draw[dashed, red!70!black, thick] (axis cs:0,0.1) -- (axis cs:0,3.8);
\node[teal!75!black] at (axis cs:0.55,1.85) {Área $=2$};
\end{axis}
\end{tikzpicture}
\end{center}

\textbf{Paso 2: Elemento diferencial.}
\[
  dA = x^{-1/2}\,dx.
\]

\textbf{Paso 3: Límites de integración.}
El punto problemático es $x=0$. Introducimos $t\to 0^+$ e integramos sobre $[t,1]$.

\textbf{Paso 4: Evaluación.}
\begin{align*}
\int_0^1 x^{-1/2}\,dx
  &= \lim_{t\to 0^+}\int_t^1 x^{-1/2}\,dx
   = \lim_{t\to 0^+}\Bigl[2\sqrt{x}\Bigr]_t^1
   = \lim_{t\to 0^+}(2-2\sqrt{t}) = 2.
\end{align*}
\[
  \boxed{\int_0^1\frac{dx}{\sqrt{x}} = 2 \quad (\text{converge}).}
\]
\end{myproof}

% ------------------------------------------------------------
\begin{example}[Discontinuidad en un punto interior]
Evalúe, si converge,
$\displaystyle\int_{0}^{3}\frac{dx}{(x-1)^{2/3}}$.
\end{example}
\begin{myproof}
\textbf{Paso 1: Análisis gráfico y región.}

El denominador se anula en $x=1\in(0,3)$; allí $f(x)\to+\infty$. La función es positiva a ambos
lados de $x=1$ y exhibe una \textbf{cúspide vertical asintótica}. Es \emph{obligatorio} separar la
integral en $[0,1)$ y $(1,3]$; de lo contrario se aplica el TFC a una función discontinua y se
obtiene un valor erróneo.

\begin{center}
\begin{tikzpicture}
\begin{axis}[
    width=0.62\textwidth, height=0.39\textheight,
    axis lines=middle,
    xlabel={$x$}, ylabel={$y$},
    xmin=-0.2, xmax=3.3, ymin=-0.28, ymax=5.1,
    restrict y to domain=0:5.1,
    samples=250,
    xtick={0,1,2,3}, ytick={0,1,2,3,4}
]
%--- Rama izquierda: x en [0, 0.94]
%    Usamos exp(-2/3*ln(|x-1|)) para evitar logaritmo de negativo
\addplot[name path=fL, thick, violet!80!black, domain=0:0.92]
    {exp(-2/3 * ln(1-x))};
%--- Rama derecha: x en [1.08, 3]
\addplot[name path=fR, thick, violet!80!black, domain=1.08:3.1]
    {exp(-2/3 * ln(x-1))};
%--- Rellenos
\path[name path=ejeL] (axis cs:0,0) -- (axis cs:0.92,0);
\path[name path=ejeR] (axis cs:1.08,0) -- (axis cs:3.1,0);
\addplot[violet!18, opacity=0.88] fill between[of=fL and ejeL];
\addplot[violet!18, opacity=0.88] fill between[of=fR and ejeR];
%--- Asíntota vertical
\draw[dashed, red!70!black, thick]
    (axis cs:1,-0.2) -- (axis cs:1,4.8)
    node[above] {$x=1$};
\node[violet!80!black] at (axis cs:0.38,2.0) {\small Izq.};
\node[violet!80!black] at (axis cs:2.0,1.5) {\small Der.};
\end{axis}
\end{tikzpicture}
\end{center}

\textbf{Paso 2: Elemento diferencial.}
\[
  dA = (x-1)^{-2/3}\,dx.
\]

\textbf{Paso 3: Límites de integración.}
Usando $t\to 1^-$ y $s\to 1^+$:
\[
  \int_0^3(x-1)^{-2/3}\,dx
  = \lim_{t\to 1^-}\int_0^t(x-1)^{-2/3}\,dx
  + \lim_{s\to 1^+}\int_s^3(x-1)^{-2/3}\,dx.
\]

\textbf{Paso 4: Evaluación.}
Recordando $\int(x-1)^{-2/3}\,dx = 3(x-1)^{1/3}+C$ (con $(x-1)^{1/3}$ para $x<1$ interpretado
como $-|x-1|^{1/3}$):
\begin{align*}
\lim_{t\to 1^-}\Bigl[3(x-1)^{1/3}\Bigr]_0^t
  &= \lim_{t\to 1^-}\bigl(3(t-1)^{1/3} - 3(-1)^{1/3}\bigr)
   = 0 - (-3) = 3. \\[4pt]
\lim_{s\to 1^+}\Bigl[3(x-1)^{1/3}\Bigr]_s^3
  &= \lim_{s\to 1^+}\bigl(3\cdot 2^{1/3} - 3(s-1)^{1/3}\bigr)
   = 3\sqrt[3]{2} - 0 = 3\sqrt[3]{2}.
\end{align*}
Ambos límites son finitos, por lo tanto la integral \textbf{converge} y:
\[
  \boxed{\int_0^3\frac{dx}{(x-1)^{2/3}} = 3 + 3\sqrt[3]{2} \approx 6.78.}
\]
\end{myproof}

% ============================================================
\section{La paradoja de la trompeta de Gabriel}
% ============================================================

Una de las aplicaciones más fascinantes del cálculo integral surge al combinar sólidos de revolución
con integrales impropias. La \textbf{trompeta de Gabriel} se genera al rotar la curva $y=1/x$ para
$x\geq 1$ alrededor del eje $x$. Esta figura desafía la intuición geométrica clásica: su volumen es
\textbf{finito}, pero su área superficial es \textbf{infinita}.

\begin{rem}[Resolución de la paradoja]
La aparente contradicción se resuelve analizando la velocidad de decaimiento de cada integrando:
\begin{itemize}
    \item \textbf{Volumen:} el integrando es $\pi/x^2$, con $p=2>1$, por lo que
          $\int_1^\infty x^{-2}\,dx$ converge.
    \item \textbf{Superficie:} el integrando es esencialmente $2\pi/x$, con $p=1$, por lo que
          $\int_1^\infty x^{-1}\,dx$ diverge.
\end{itemize}
La paradoja surge de intentar aplicar intuición física de «grosor uniforme» a un objeto matemático
abstracto. En el modelo teórico, «llenar» la trompeta equivale a adelgazar la pintura exactamente
como $1/x^2$ a medida que $x\to\infty$, lo cual es perfectamente consistente.
\end{rem}

% ------------------------------------------------------------
\begin{example}[Volumen finito y área superficial infinita de la trompeta]
Demuestre que el volumen de la trompeta de Gabriel es $V=\pi$ (finito), pero que su área
superficial es infinita.
\end{example}
\begin{myproof}
\textbf{Paso 1: Análisis gráfico y región.}

La curva $y=1/x$ rota alrededor del eje $x$ para $x\geq 1$. Los discos perpendiculares al eje
tienen radio $R(x)=1/x$. La región generatriz es el área entre la curva, el eje $x$ y la recta
$x=1$.

\begin{center}
\begin{tikzpicture}
\begin{axis}[
    width=0.65\textwidth, height=0.40\textheight,
    axis lines=middle,
    xlabel={$x$}, ylabel={$y$},
    xmin=0.5, xmax=5.1, ymin=-0.12, ymax=1.22,
    domain=1:4.9, samples=130,
    xtick={1,2,3,4}, ytick={0,0.25,0.5,0.75,1}
]
\addplot[name path=f, thick, blue!80!black] {1/x};
\addlegendentry{$y=1/x$}
\path[name path=eje] (axis cs:1,0) -- (axis cs:4.9,0);
\addplot[blue!18, opacity=0.85]
    fill between[of=f and eje, soft clip={domain=1:4.9}];
\draw[dashed, red!70!black, thick]
    (axis cs:1,0) -- (axis cs:1,1) node[left] {$(1,1)$};
%--- Indicador de radio en x=2.5
\draw[gray!70!black, <->, thick]
    (axis cs:2.5,0) -- (axis cs:2.5,{1/2.5})
    node[midway, right] {\small $R(x)=\tfrac{1}{x}$};
\node[blue!60!black] at (axis cs:3.4,0.18) {Región generatriz};
\end{axis}
\end{tikzpicture}
\end{center}

\textbf{Paso 2: Elemento diferencial.}
\begin{itemize}
    \item \textbf{Volumen} (método del disco):
          $dV = \pi[R(x)]^2\,dx = \pi x^{-2}\,dx$.
    \item \textbf{Superficie} (fórmula de área de superficie de revolución, con $f'(x)=-x^{-2}$):
          $dS = 2\pi f(x)\sqrt{1+[f'(x)]^2}\,dx
              = \dfrac{2\pi}{x}\sqrt{1+\dfrac{1}{x^4}}\,dx$.
\end{itemize}

\textbf{Paso 3: Límites de integración.}
Ambas integrales se evalúan sobre $[1,\infty)$.

\textbf{Paso 4: Evaluación.}

\textit{Volumen:}
\begin{align*}
V &= \pi\int_1^\infty x^{-2}\,dx
   = \pi\lim_{b\to\infty}\Bigl[-x^{-1}\Bigr]_1^b
   = \pi\lim_{b\to\infty}\Bigl(1-\frac{1}{b}\Bigr)
   = \pi.
\end{align*}

\textit{Superficie} (criterio de comparación directa, usando que
$\sqrt{1+x^{-4}}>1$ para todo $x\geq 1$):
\begin{align*}
A = 2\pi\int_1^\infty\frac{\sqrt{1+x^{-4}}}{x}\,dx
  &> 2\pi\int_1^\infty\frac{dx}{x}
   = 2\pi\lim_{b\to\infty}[\ln x]_1^b
   = +\infty.
\end{align*}
Por el criterio de comparación directa, la integral del área \textbf{diverge}. Concluimos:
\[
  \boxed{V = \pi \quad \text{(finito)} \qquad \text{y} \qquad A = +\infty \quad \text{(diverge).}}
\]
\end{myproof}

% ============================================================
%  PROBLEMAS PROPUESTOS — INTEGRALES IMPROPIAS
% ============================================================
\section{Problemas Propuestos}

Recuerde aplicar los criterios de convergencia y las definiciones de límite
estudiadas en clase para justificar la existencia de cada valor en las
integrales impropias. En todos los casos, el desarrollo analítico completo
constituye parte esencial de la solución.


\begin{prob}
Determine si las siguientes integrales impropias de Tipo~I (con al menos
un límite de integración infinito) convergen o divergen. En caso de
convergencia, calcule su valor exacto.
\begin{enumerate}
    \item $\displaystyle \int_{0}^{\infty} \frac{dx}{1 + x^2}$
    \item $\displaystyle \int_{-\infty}^{0} e^{3x} \, dx$
    \item $\displaystyle \int_{0}^{\infty} x e^{-x} \, dx$
    \item $\displaystyle \int_{1}^{\infty} \frac{\ln x}{x} \, dx$
    \item $\displaystyle \int_{-\infty}^{\infty} \frac{x}{x^2 + 1} \, dx$
    \item $\displaystyle \int_{0}^{\infty} \frac{e^{x}}{e^{2x} + 1} \, dx$
\end{enumerate}
\end{prob}

\begin{prob}
Evalúe cada una de las siguientes integrales impropias, o demuestre que
divergen.
\begin{enumerate}
    \item $\displaystyle \int_{1}^{\infty} \frac{1}{x^3} \, dx$
    \item $\displaystyle \int_{0}^{\infty} e^{-2x} \, dx$
    \item $\displaystyle \int_{-\infty}^{\infty} \frac{1}{x^2 + 4} \, dx$
    \item $\displaystyle \int_{0}^{\infty} x^2 e^{-x} \, dx$
    \item $\displaystyle \int_{2}^{\infty} \frac{1}{x \ln x} \, dx$
    \item $\displaystyle \int_{0}^{\infty} \cos x \, dx$
\end{enumerate}
\end{prob}

\begin{prob}
Determine si la siguiente integral converge y, de ser así, encuentre su
valor exacto:
\[
    \int_{-\infty}^{\infty} \frac{e^x}{1 + e^{2x}} \, dx.
\]
\textit{Sugerencia: divida la integral en los intervalos $(-\infty, 0]$ y
$[0, +\infty)$, y utilice la sustitución $u = e^x$ en cada parte.}
\end{prob}

\begin{prob}
Analice la convergencia de las siguientes integrales impropias de Tipo~II,
cuyo integrando presenta una discontinuidad infinita en uno o ambos
extremos del intervalo de integración. Si la integral converge, determine
su valor exacto.
\begin{enumerate}
    \item $\displaystyle \int_{0}^{16} \frac{1}{\sqrt{x}} \, dx$
    \item $\displaystyle \int_{0}^{1} \ln x \, dx$
    \item $\displaystyle \int_{0}^{2} \frac{dx}{\sqrt{4 - x^2}}$
    \item $\displaystyle \int_{0}^{\pi/2} \tan x \, dx$
    \item $\displaystyle \int_{1}^{e} \frac{dx}{x\sqrt{\ln x}}$
    \item $\displaystyle \int_{0}^{1} \frac{x}{\sqrt{1 - x^2}} \, dx$
\end{enumerate}
\end{prob}

\begin{prob}
Evalúe las siguientes integrales impropias, en las cuales el integrando
presenta una discontinuidad infinita en un punto interior $c$ del intervalo
$[a, b]$ (con $a < c < b$), o demuestre que divergen.
\begin{enumerate}
    \item $\displaystyle \int_{0}^{3} \frac{dx}{(x-1)^{2/3}}$
    \item $\displaystyle \int_{0}^{2} \frac{dx}{x - 1}$
    \item $\displaystyle \int_{-3}^{3} \frac{x}{\sqrt{9 - x^2}} \, dx$
    \item $\displaystyle \int_{-1}^{2} \frac{1}{x^3} \, dx$
    \item $\displaystyle \int_{-1}^{1} \frac{1}{\sqrt{|x|}} \, dx$
\end{enumerate}
\end{prob}


\begin{prob}
La siguiente integral es impropia por dos razones simultáneas: el
integrando presenta una singularidad en $x = 0$ y el intervalo de
integración no está acotado. Descomponga el intervalo en $x = 1$,
justifique la convergencia de cada parte por separado mediante la
evaluación del límite correspondiente y, si ambas subintegrales convergen,
determine el valor exacto de la integral completa:
\[
    \int_{0}^{\infty} \frac{dx}{\sqrt{x}\,(1 + x)}.
\]
\end{prob}

\begin{prob}
Analice la convergencia de la siguiente integral, que es impropia tanto
por la singularidad del integrando en $x = 0$ como por el límite superior
de integración infinito. Descomponga el intervalo en $x = 1$, estudie
cada parte por separado y, si ambas convergen, calcule el valor exacto:
\[
    \int_{0}^{\infty} x^{-1/3} e^{-x} \, dx.
\]
\end{prob}

\begin{prob}
Considere la integral
\[
    \int_{0}^{\infty} \frac{\ln x}{1 + x^2} \, dx.
\]
\begin{enumerate}
    \item Identifique con precisión las dos fuentes de impropie\-dad de esta
          integral.
    \item Descomponga el dominio en $x = 1$ y estudie la convergencia de
          cada parte por separado.
    \item Utilice la sustitución $x = 1/t$ en la integral sobre $(0,1]$
          para demostrar que esa parte es el negativo de la integral sobre
          $[1, \infty)$ y concluya el valor de la integral completa.
\end{enumerate}
\end{prob}

\begin{prob}
Considere la integral de tipo $p$ definida sobre un intervalo no acotado.
Demuestre formalmente, mediante la evaluación explícita del límite
correspondiente, que la integral
\[
    \int_{1}^{\infty} \frac{1}{x^p} \, dx
\]
\begin{enumerate}
    \item diverge si $p \leq 1$, y
    \item converge si $p > 1$; en este caso, determine su valor exacto
          en función del parámetro $p$.
\end{enumerate}
\end{prob}

\begin{prob}
De manera análoga al problema anterior, demuestre formalmente que la
integral
\[
    \int_{0}^{1} \frac{1}{x^p} \, dx
\]
\begin{enumerate}
    \item converge si $p < 1$; en este caso, calcule su valor exacto en
          función de $p$, y
    \item diverge si $p \geq 1$.
\end{enumerate}
\end{prob}

\begin{prob}
Utilice el Teorema de Comparación Directa o el Teorema de Comparación en
el Límite para determinar la convergencia o divergencia de cada integral.
No es necesario calcular el valor exacto; es suficiente identificar con
precisión la función de comparación y justificar rigurosamente la
conclusión.
\begin{enumerate}
    \item $\displaystyle \int_{1}^{\infty} e^{-x^2} \, dx$
    \item $\displaystyle \int_{1}^{\infty} \frac{\sin^2 x}{x^2} \, dx$
    \item $\displaystyle \int_{2}^{\infty} \frac{dx}{\sqrt{x^4 - 1}}$
    \item $\displaystyle \int_{1}^{\infty} \frac{x + 1}{\sqrt{x^4 + x}} \, dx$
    \item $\displaystyle \int_{1}^{\infty} \frac{\arctan x}{1 + x^2} \, dx$
\end{enumerate}
\end{prob}

\begin{prob}
Para cada una de las siguientes integrales, determine si converge o
diverge. Justifique su respuesta mediante un criterio de comparación
apropiado, identificando explícitamente la función auxiliar utilizada y
verificando las hipótesis del criterio elegido.
\begin{enumerate}
    \item $\displaystyle \int_{0}^{1} \frac{\sin x}{x^{3/2}} \, dx$
    \item $\displaystyle \int_{1}^{\infty} \frac{x^2 + 1}{x^4 + x^2 + 1} \, dx$
    \item $\displaystyle \int_{1}^{\infty} \frac{\ln x}{x^{3/2}} \, dx$
    \item $\displaystyle \int_{0}^{1} \frac{e^x - 1}{x^{1/2}} \, dx$
    \item $\displaystyle \int_{2}^{\infty} \frac{1}{\ln x} \, dx$
\end{enumerate}
\end{prob}

\begin{prob}[Función Gamma]
La función Gamma se define para $n > 0$ como:
\[
    \Gamma(n) = \int_{0}^{\infty} x^{n-1} e^{-x} \, dx.
\]
\begin{enumerate}
    \item Use integración por partes para demostrar la relación de
          recurrencia $\Gamma(n+1) = n\,\Gamma(n)$.
    \item Deduzca que si $n$ es un entero positivo, entonces
          $\Gamma(n) = (n-1)!$
    \item Calcule explícitamente $\Gamma(1)$, $\Gamma(2)$ y $\Gamma(3)$.
\end{enumerate}
\end{prob}

\begin{prob}[Probabilidad]
Sea $X$ una variable aleatoria continua con función de densidad de
probabilidad
\[
    f(x) =
    \begin{cases}
        0.01\,e^{-0.01x} & \text{si } x \geq 0, \\
        0 & \text{si } x < 0.
    \end{cases}
\]
\begin{enumerate}
    \item Verifique que $f$ es efectivamente una función de densidad de
          probabilidad, es decir, que $\displaystyle\int_{-\infty}^{\infty} f(x)\,dx = 1$.
    \item Calcule $P(X > 20)$ mediante la evaluación de la integral
          impropia pertinente.
    \item Determine el valor esperado $E[X] = \displaystyle\int_{0}^{\infty} x\,f(x)\,dx$.
\end{enumerate}
\end{prob}

\begin{prob}[Geometría Analítica]
El piloto de un automóvil de carreras recorre una pista elíptica modelada
por la ecuación
\[
    \frac{x^2}{400} + \frac{y^2}{100} = 1.
\]
Al perder el control en el punto $(16,\,6)$, el vehículo continúa en línea
recta sobre la recta tangente a la elipse en ese punto.
\begin{enumerate}
    \item Verifique que el punto $(16,\,6)$ pertenece a la elipse.
    \item Determine la ecuación de la recta tangente a la elipse en dicho
          punto.
    \item Si el automóvil impacta contra un árbol ubicado en la coordenada
          $(14,\, k)$, determine el valor exacto de la constante $k$.
\end{enumerate}
\end{prob}

\begin{prob}[Longitud de arco]
La longitud de arco de una curva $y = f(x)$ sobre el intervalo $[a, b]$
está dada por la integral
\[
    L = \int_{a}^{b} \sqrt{1 + \left[f'(x)\right]^2} \, dx.
\]
\begin{enumerate}
    \item Muestre que la longitud de la curva $y = \sqrt{1 - x^2}$ sobre
          $[0,1)$ está dada por una integral impropia de Tipo~II y evalúe
          dicha integral.
    \item Interprete geométricamente el resultado obtenido.
\end{enumerate}
\end{prob}

\begin{prob}[Física — Trabajo gravitacional]
La fuerza de atracción gravitacional entre la Tierra (masa $M$) y un objeto
de masa $m$ situado a una distancia $r$ de su centro es
\[
    F(r) = \frac{GMm}{r^2},
\]
donde $G$ es la constante de gravitación universal y $R$ denota el radio
de la Tierra.
\begin{enumerate}
    \item Plantee la integral impropia que representa el trabajo necesario
          para trasladar el objeto desde la superficie terrestre ($r = R$)
          hasta el infinito.
    \item Evalúe dicha integral y exprese el resultado en términos de $G$,
          $M$, $m$ y $R$.
    \item Explique por qué la convergencia de esta integral implica que
          existe una velocidad de escape finita.
\end{enumerate}
\end{prob}

